% =================================================================
% doc-class.tex — Document class & ABNT formatting (NBR 14724/6024)
% Citations are handled by pandoc + CSL (csl/abnt.csl), not biblatex.
% =================================================================
\documentclass[a4paper,12pt]{article}

% Encoding & language
\usepackage[utf8]{inputenc}
\usepackage[T1]{fontenc}

% Page geometry: ABNT NBR 14724 (3cm top/left, 2cm bottom/right)
\usepackage[lmargin=3cm,tmargin=3cm,rmargin=2cm,bmargin=2cm]{geometry}

% Page numbering top-right (ABNT)
\usepackage{fancyhdr}
\fancyhf{}
\fancyhead[R]{\thepage}
\renewcommand{\headrulewidth}{0pt}

% Line spacing: 1.5 between lines (ABNT)
\usepackage[onehalfspacing]{setspace}

% Justified text (ABNT) + first-line indent on every paragraph
\usepackage[document]{ragged2e}
\usepackage{indentfirst}
\usepackage[skip=5pt, indent=15pt]{parskip}

% Layout helpers
\usepackage{graphicx, xcolor}
\usepackage{enumerate, multirow, multicol}
\usepackage{adjustbox, changepage}
\usepackage{float}
\usepackage{pdflscape, pdfpages}

% Math
\usepackage{amsfonts, amsthm, amsmath, amssymb, mathtools}

% TikZ (optional; comment out if not used)
\usepackage{tikz}

% SVG include support (requires Inkscape on the system)
\usepackage[notransparent]{svg}

% Section title formatting (ABNT NBR 6024)
\usepackage{titlesec}
\titleformat{\section}      {\normalsize\bfseries}{\thesection.}      {1em}{}
\titleformat{\subsection}   {\normalsize}         {\thesubsection.}   {1em}{}
\titleformat{\subsubsection}{\normalsize}         {\thesubsubsection.}{1em}{}

% Conditional helper
\usepackage{ifthen}

% ----------------------------------------------------------------
% ABNT NBR 10520 — long quote (citação direta com mais de 3 linhas)
% ----------------------------------------------------------------
\newlength{\ABNTEXcitacaorecuo}
\setlength{\ABNTEXcitacaorecuo}{4cm}
\newcommand{\ABNTEXfontereduzida}{\footnotesize}
\newenvironment*{citacao}[1][default]{%
   \list{}%
   \ABNTEXfontereduzida%
   \addtolength{\leftskip}{\ABNTEXcitacaorecuo}%
   \item[]%
   \singlespacing
   \ifthenelse{\not\equal{#1}{default}}{\itshape\selectlanguage{#1}}{}%
}{%
   \endlist%
}

% ----------------------------------------------------------------
% Convenience macros for figures and quadros (ABNT)
%   \fig{caption}{label}{path}{source}
% ----------------------------------------------------------------
\newcommand{\fig}[4]{%
  \begin{figure}[H]
    \centering
    \caption{#1}
    \label{#2}
    \includegraphics{#3}
    \vspace{0.5cm}
    \begin{footnotesize}Fonte: #4\end{footnotesize}
  \end{figure}
}

\newcommand{\quadro}[4]{%
  \renewcommand{\figurename}{Quadro}
  \setcounter{figure}{0}
  \begin{figure}[H]
    \captionsetup{list=no}
    \centering
    \caption{#1}
    \label{#2}
    \includegraphics{#3}
    \vspace{0.5cm}
    \begin{footnotesize}Fonte: #4\end{footnotesize}
  \end{figure}
}

% Footnotesize helper for tables
\newcommand{\FontTable}[1]{\vspace{-\baselineskip}\begin{footnotesize}#1\end{footnotesize}}

% Hidden text (used in pre-textual blocks for vertical alignment)
\newenvironment{invisible}{\color{white}}{}

% Common math macros
\newcommand{\PR}[1]{\ensuremath{\left[#1\right]}}
\newcommand{\PC}[1]{\ensuremath{\left(#1\right)}}
\newcommand{\chav}[1]{\ensuremath{\left\{#1\right\}}}
\newcommand{\floor}[1]{\ensuremath{\left\lfloor #1\right\rfloor}}
